\textbf{Construcción de la matriz y producto por filas.} Se convierte cada ángulo a radianes y se escribe $c_\phi=\cos\phi$, $s_\phi=\sin\phi$, y análogamente para los otros ángulos. Multiplicar las rotaciones pasivas en orden 3--2--1 da:\[\fitmath{C_{bn}=\begin{bmatrix}c_\theta c_\psi&c_\theta s_\psi&-s_\theta\\s_\phi s_\theta c_\psi-c_\phi s_\psi&s_\phi s_\theta s_\psi+c_\phi c_\psi&s_\phi c_\theta\\c_\phi s_\theta c_\psi+s_\phi s_\psi&c_\phi s_\theta s_\psi-s_\phi c_\psi&c_\phi c_\theta\end{bmatrix}}\]\[\fitmath{u=(0.837607217)(96)+(0.523395079)(8)+(-0.156434465)(-5)\approx85.37963\,\mathrm{m/s}}\]\[\fitmath{v=(-0.462987754)(96)+(0.832158446)(8)+(0.305212482)(-5)\approx-39.31562\,\mathrm{m/s}}\]\[\fitmath{w=(0.289924973)(96)+(-0.183220936)(8)+(0.939347432)(-5)\approx21.67029\,\mathrm{m/s}}\]La inversa es la transpuesta: multiplicar $C_{bn}^T$ por el vector obtenido recupera las tres componentes NED. Las normas coinciden: $V=96.46243\,\mathrm{m/s}$. Para incidencia se usa $\alpha=\operatorname{atan2}(w,u)$; el argumento de $w/u$ por sí solo pierde el cuadrante.
